% META: DONE.1

\section{Реализация нейросетевой части}

\subsection{Обоснование выбора MeshCNN}
Для задач сегментации трёхмерных сеток наиболее подходящей архитектурой является MeshCNN \cite{meshcnn2019}. В отличие от PointNet или DGCNN, работающих с облаками точек, MeshCNN оперирует непосредственно на рёбрах треугольной сетки, сохраняя топологию. Это критически важно для распознавания таких инженерных признаков, как скругления (плавное изменение кривизны) и фаски (плоские наклонные грани). Кроме того, MeshCNN использует инвариантные признаки (углы между гранями, синусы двугранных углов), что делает её устойчивой к разным способам триангуляции одной и той же геометрии.

\subsection{Архитектура оригинальной MeshCNN}
Каждая модель (классификатор или сегментатор) построена на базе стандартной MeshCNN с незначительными модификациями. Основные компоненты:

\begin{itemize}
    \item \textbf{Входной слой}: Для всех экспериментов фиксировано число рёбер на входе сети: $N_{\text{edges}} = 2250$. Это значение выбрано как компромисс между сохранением геометрических деталей и вычислительной эффективностью: модели с меньшим числом рёбер (например, 750) теряют мелкие признаки скруглений и фасок, а с большим - не помещаются в память графического процессора при размере батча 4-8. Каждое ребро характеризуется 5 инвариантными признаками: 
        \begin{itemize}
            \item двугранный угол между инцидентными гранями;
            \item два внутренних угла (при вершинах ребра в каждом из двух прилегающих треугольников);
            \item два отношения длины ребра к высоте (для каждого треугольника - длина ребра, делённая на длину перпендикуляра к противолежащей вершине).
        \end{itemize}
        Эти признаки инвариантны к повороту, переносу и равномерному масштабированию.
    \item \textbf{Свёрточные блоки}: последовательность из нескольких остаточных свёрточных слоёв (ResConv) с ядром 5 (учитываются ребро и его 4 соседа). Количество каналов различалось среди экспериментов, но в итоге для сегментации было выбрано: 32 -> 32 -> 64 -> 64 -> 128, а для классификации: 32 -> 64 -> 128. Причём после каждого "основного" блока дополнительно добавляются 1-3 аналогичных свёрточных блоков, соединённых последовательно, с skip\_connections. После каждой свёртки применяется групповая нормализация.
    \item \textbf{Пулинг}: после каждого свёрточного блока выполняется пулинг на основе коллапса рёбер (edge collapse), последовательно уменьшая число рёбер. Для сегментации выбрано 2400 -> 1800 -> 1200 -> 600 (после последней свёртки не добавляется пуллинг), а для классификации: 2400 -> 1800 -> 1200. Сеть сама учится, какие рёбра сворачивать, ориентируясь на минимизацию функции потерь, таким образом получая возможность выделять фичи и оставлять только нужные фрагменты модели, сохраняя как локальные так и глобальные характеристики.
    \item \textbf{Классификационная голова}: после последнего пулинга выполняется глобальный средний пулинг (Global Average Pooling), затем полносвязный слой с выходом 2 (вероятность наличия операции).
    \item \textbf{Сегментационная голова}: симметричная архитектура с анпулингом (mesh unpooling), восстанавливающим исходное число рёбер.
\end{itemize}

\subsection{Параметры оптимизации}
Использовался оптимизатор Adam с параметрами $\beta_1 = 0.9$, $\beta_2 = 0.999$, начальная скорость обучения $lr = 0.001$ для сегментации и $0.002$ для классификации (подобрано эмпирически). Политика изменения скорости обучения - \texttt{lambda}: 
\[
lr_{\text{epoch}} = lr_{\text{init}} \cdot \left(1 - \max\left(0, \frac{\text{epoch} + 1 - \text{niter}}{\text{niter\_decay} + 1}\right)\right),
\]
где $\text{niter}=10$ - число эпох с постоянной скоростью, $\text{niter\_decay}=10\ldots40$ - число эпох линейного затухания до нуля. Размер батча составлял $8$ для классификации и $4$ для сегментации (ограничение видеопамяти).

\subsection{Функции потерь и оптимизация}
Для классификации используется бинарная кросс-энтропия:
\[
\mathcal{L}_{\text{cls}} = -\frac{1}{N}\sum_{i=1}^{N}\left[y_i\log p_i + (1-y_i)\log(1-p_i)\right],
\]
где $y_i$ - истинная метка модели (0 - операция отсутствует, 1 - присутствует), $p_i$ - предсказанная вероятность.

Для сегментации применяется взвешенная кросс-энтропия, чтобы скомпенсировать дисбаланс классов (фоновых рёбер обычно значительно больше, чем рёбер, принадлежащих операции):
\[
\mathcal{L}_{\text{seg}} = -\frac{1}{M}\sum_{j=1}^{M}\sum_{c\in\{0,1\}} w_c \cdot y_{j,c} \log p_{j,c},
\]
где $w_0$ и $w_1$ выбираются обратно пропорционально частоте классов в обучающей выборке.

\subsection{Аугментация данных}
Для повышения обобщающей способности применялись следующие методы аугментации (включены в \texttt{train.sh}):
\begin{itemize}
    \item \textbf{Скольжение вершин} (\texttt{--slide\_verts 0.2}) - 20\% вершин случайно смещаются вдоль поверхности сетки на расстояние до 30\% длины инцидентного ребра. Это имитирует незначительные изменения триангуляции.
    \item \textbf{Переворот рёбер} (\texttt{--flip\_edges 0.2}) - 20\% рёбер случайным образом переворачиваются (flip edge), что изменяет топологию сетки, сохраняя геометрию.
    \item \textbf{Масштабирование вершин} (\texttt{--scale\_verts}) - неоднородное масштабирование координат вершин по каждой оси с коэффициентом из нормального распределения $\mathcal{N}(1, 0.1)$.
\end{itemize}
Количество вариантов аугментации на одну модель задавалось параметром \texttt{--num\_aug 8}, что позволило эффективно увеличить обучающую выборку в 8 раз без повторной генерации размеченных данных.

